\documentclass[UTF-8]{ctexart}
\usepackage{geometry}
\geometry{a4paper, margin=2cm}
\usepackage{graphicx}
\usepackage{longtable}
\usepackage{booktabs}
\usepackage{enumitem}

\begin{document}

\title{需求分析}
\author{}
\date{}
\maketitle
\newpage

\tableofcontents
\newpage

\section{简介}
\subsection{目的}
本文档旨在详细描述图书管理系统的功能需求、非功能需求、数据需求和实现方案，为系统开发和测试提供指导依据

\subsection{名词定义}

\begin{longtable}{ll}
\toprule
名词 & 解释 \\
\midrule
ISBN & 国际标准书号，图书的唯一标识符 \\
借阅记录 & 记录图书借阅和归还信息的数据 \\
库存 & 图书的可用数量 \\
未归还 & 已借出但尚未归还的图书状态 \\
\bottomrule
\end{longtable}

\section{总体概述}
\subsection{软件概述}
图书管理系统是一款专为图书管理场景设计的软件系统，为用户提供图书信息管理和借阅流程跟踪的数字化工具，实现图书全生命周期管理和借阅业务的规范化、智能化处理。
系统设计简洁明了，功能完整，易于使用和维护，短小精悍，适合小型图书馆、学校图书室或个人图书收藏管理使用，为图书管理工作提供高效、便捷的数字化解决方案。

\subsection{软件功能}
本项目输出产品名称为图书管理系统，是一个用于管理图书和借阅流程的软件系统，在这个系统中将实现以下特性：
\begin{enumerate}[label=\arabic{enumi}、]
\item 支持图书信息的管理，包括新增、查询、修改、删除图书。 （优先级：重要）
\item 支持图书借阅流程管理，包括借阅图书、归还图书、查看借阅记录。 （优先级：重要）
\item 支持按多种方式查询图书，包括按书名、作者、ISBN查询。（优先级：重要）
\item 支持借阅记录的多维度查询，包括所有记录、指定借阅人记录、未归还记录。 （优先级：一般）
\item 支持数据的本地存储，保存图书和借阅记录数据。 （优先级：重要）
\item 提供图形化操作界面，满足不同场景的使用需求。 （优先级：一般）
\item 支持库存管理，自动处理借阅和归还时的库存变化。（优先级：一般）
\end{enumerate}

图书管理系统不支持以下特性： 
\begin{enumerate}[label=\arabic{enumi}、]
\item 不支持用户权限管理，所有用户拥有相同的操作权限。 
\item 不支持图书分类管理，仅通过基本信息进行管理。 
\item 不支持借阅期限和逾期处罚功能。 
\item 不支持图书预约功能。 
\item 不支持数据备份和恢复功能。
\end{enumerate}

向用户提供以下操作接口：
\begin{itemize}
\item 导航栏：首页、图书管理、借阅管理。
\item 图书管理页：图书搜索、图书列表、新增/修改/删除操作。
\item 借阅管理页：借阅/归还操作、借阅记录查询。
\item 首页：系统功能介绍。
\end{itemize}

\subsection{用户特征}
\begin{enumerate}[label=\arabic{enumi}]
\item 图书管理员 ：
\begin{itemize}
\item 角色：负责图书的采购、分类、编目等管理工作。
\item 特征：了解图书分类和编目知识，具备基本的计算机操作技能。
\item 需求：需要便捷的图书信息管理功能，能够快速录入和修改图书信息。
\end{itemize}

\item 系统维护人员 ：
\begin{itemize}
\item 角色：负责系统的安装、配置和维护。
\item 特征：具备基本的计算机技术知识。
\item 需求：需要系统稳定可靠，数据备份方便，易于维护和扩展。
\end{itemize}

\item 普通用户：
\begin{itemize}
\item 角色：查询图书信息和个人借阅记录。
\item 特征：具备基本的计算机操作能力。
\item 需求：需要简单直观的查询界面，能够快速找到所需图书信息。
\end{itemize}
\end{enumerate}

\subsection{可行性分析}
本项目为软件工程课程大作业，围绕图书管理系统的开发目标，从四大核心维度论证落地可行性。

\subsubsection{技术可行性}
系统基于Web和Python 开发，核心功能仅需 Python 标准库即可实现，扩展版本采用轻量开源框架，无技术壁垒；兼容 Windows、macOS、Linux 全主流操作系统，仅需普通个人电脑即可完成开发与运行，完全适配课程所学技术栈，开发落地无风险。 

\subsubsection{操作可行性}
系统采用直观的层级菜单与标准化图形界面设计，操作流程贴合图书管理实际业务场景，配套清晰的输入指引与友好的错误提示；目标用户仅需具备基础计算机操作能力即可快速上手，无需专业技术培训，使用门槛低。 

\subsubsection{经济可行性}
项目全流程零额外成本投入，所用开发工具、框架、依赖库均为开源免费资源，无版权、授权、硬件采购、运维等相关费用。

\subsubsection{进度可行性}
系统采用模块化拆分设计，核心业务逻辑清晰，可分模块开发与测试；整体开发周期短，项目进度完全可控。 

综上，本项目在技术、操作、经济、进度维度均具备完整的落地可行性，可正式启动开发工作。

\section{系统功能性分析}
\subsection{基于场景的建模}
在本图书管理系统的分析与建模过程中，我们围绕三大核心用例模块，从actor的角度出发，构建了覆盖全业务流程的场景模型。

\subsubsection{总体用例图}
总用例图如图3.1.1所示：
\begin{figure}
图3.1.1总体用例图
    \centering
    \includegraphics[width=0.8\textwidth]{用例图.png}
\end{figure}

\subsubsection{总体泳道图}
总体泳道图如图3.1.2所示：
\begin{figure}
图3.1.2总体泳道图
    \centering
    \includegraphics[width=0.8\textwidth]{泳道图.png}
\end{figure}

\subsubsection{总体活动图}
总体活动图如图3.1.3所示：
\begin{figure}
图3.1.3总体活动图
    \centering
\end{figure}

\subsubsection{子模块分析}
\paragraph{查询图书用例分析及建模}
\leavevmode 
\begin{table}[htbp]
\centering
\begin{tabular}{|p{3cm}|p{12cm}|}
\hline
\textbf{用例描述} & 用户选择查询维度，输入关键词检索图书信息，系统返回匹配的图书列表，是所有图书操作的前置基础能力 \\
\hline
\textbf{参与者} & 图书管理员 \\
\hline
\textbf{前置条件} & 1. 用户已进入系统，图书管理模块可正常访问；2. 系统图书数据已完成加载 \\
\hline
\textbf{后置条件} & 1. 系统展示匹配的图书列表，无匹配结果时给出明确提示；2. 不修改系统任何存储数据 \\
\hline
\textbf{基本事件流} & 1. 用户进入图书管理模块，进入查询功能页面2. 用户选择查询方式3. 用户输入查询关键词，提交查询请求4. 系统根据查询维度与关键词，遍历图书数据执行匹配5. 系统渲染匹配的图书列表，完成查询操作 \\
\hline
\textbf{意外事件流} & 1. 若关键词为空，系统提示 ，终止查询2. 若无匹配的图书数据，系统提示 ，展示空列表3. 若系统数据未加载完成，系统提示 ，终止查询 \\
\hline
\textbf{活动图} & \\
\hline
\end{tabular}
\caption{查询图书用例分析}
\end{table}

\paragraph{新增图书用例分析及建模}
\leavevmode 
\begin{table}[htbp]
\centering
\begin{tabular}{|p{3cm}|p{12cm}|}
\hline
\textbf{用例描述} & 图书管理员录入图书基础信息，系统校验 ISBN 唯一性后，完成图书入库与信息存储 \\
\hline
\textbf{参与者} & 图书管理员 \\
\hline
\textbf{前置条件} & 1. 管理员已进入图书管理模块，拥有图书操作权限；2. 系统数据已正常加载完成 \\
\hline
\textbf{后置条件} & 1. 图书信息成功写入系统本地存储；2. 图书库存同步更新；3. 图书列表自动刷新，展示新增数据 \\
\hline
\textbf{基本事件流} & 1. 管理员点击“新增图书”入口，进入信息录入页面2. 管理员录入图书必填信息3. 系统校验 ISBN 是否已存在于图书库中4. 校验通过，系统将图书信息写入本地存储5. 系统提示，刷新图书列表，完成操作 \\
\hline
\textbf{意外事件流} & 1. 若 ISBN 已存在，返回信息，终止操作，保留已录入信息2. 若必填项为空，返回提示，终止提交3. 若库存输入为负数，返回提示，终止提交4. 若数据存储失败，返回提示，回滚操作 \\
\hline
\textbf{活动图} & \\
\hline
\end{tabular}
\caption{新增图书用例分析}
\end{table}

\paragraph{修改图书信息用例分析及建模}
\leavevmode 
\begin{table}[htbp]
\centering
\begin{tabular}{|p{3cm}|p{12cm}|}
\hline
\textbf{用例描述} & 图书管理员对已入库的图书信息进行编辑更新，系统校验后完成数据修改，必须先通过查询定位目标图书 \\
\hline
\textbf{参与者} & 图书管理员 \\
\hline
\textbf{前置条件} & 1. 管理员已通过查询功能定位到目标图书；2. 系统数据已正常加载完成 \\
\hline
\textbf{后置条件} & 1. 目标图书的信息在存储中完成更新；2. 图书列表同步刷新，展示最新信息 \\
\hline
\textbf{基本事件流} & 1. 管理员在图书查询结果中，进行修改2. 系统回显该图书的当前信息，进入编辑页面3. 管理员修改图书信息4. 管理员提交修改，系统校验修改后信息的合法性5. 校验通过，系统更新存储中的图书数据6. 系统提示，刷新图书列表，完成操作 \\
\hline
\textbf{意外事件流} & 1. 若修改后的库存为负数，系统提示 ，终止提交2. 若必填项修改为空，系统提示 ，终止提交3. 若数据更新失败，系统提示 ，回滚数据，保留原信息 \\
\hline
\textbf{活动图} & \\
\hline
\end{tabular}
\caption{修改图书信息用例分析}
\end{table}

\paragraph{删除图书用例及建模}
\leavevmode 
\begin{table}[htbp]
\centering
\begin{tabular}{|p{3cm}|p{12cm}|}
\hline
\textbf{用例描述} & 图书管理员对目标图书执行下架删除操作，系统校验无未归还借阅记录后完成删除，必须先通过查询定位目标图书 \\
\hline
\textbf{参与者} & 图书管理员 \\
\hline
\textbf{前置条件} & 1. 管理员已通过查询功能定位到目标图书；2. 系统借阅记录数据已正常加载 \\
\hline
\textbf{后置条件} & 1. 目标图书从系统存储中永久移除；2. 图书列表同步刷新，移除该条数据 \\
\hline
\textbf{基本事件流} & 1. 管理员在图书查询结果中，进行删除 2. 系统弹出二次确认弹窗，提示管理员确认删除操作3. 管理员确认删除，系统校验该图书是否存在未归还的借阅记录4. 校验无未归还记录，系统从存储中删除该图书数据5. 系统提示，刷新图书列表，完成操作 \\
\hline
\textbf{意外事件流} & 1. 若该图书存在未归还的借阅记录，系统提示 "该图书有未归还的借阅记录"，终止操作2. 若管理员取消确认，系统终止删除操作，保留原数据3. 若数据删除失败，系统提示 ，保留原数据 \\
\hline
\textbf{活动图} & \\
\hline
\end{tabular}
\caption{删除图书用例分析}
\end{table}


\paragraph{借阅图书用例及建模}
\leavevmode 
\begin{table}[htbp]
\centering
\begin{tabular}{|p{3cm}|p{12cm}|}
\hline
\textbf{用例描述} & 图书管理员为借阅人办理图书借阅手续，系统校验库存、重复借阅情况后，完成借阅操作，同步扣减图书库存、生成借阅记录 \\
\hline
\textbf{参与者} & 图书管理员 \\
\hline
\textbf{前置条件} & 1. 管理员已进入借阅管理模块；2. 目标图书已在系统中存在，可通过 ISBN 定位；3. 系统图书与借阅记录数据已正常加载 \\
\hline
\textbf{后置条件} & 1. 系统生成有效借阅记录，状态标记为 “未归还”；2. 对应图书的库存数量扣减 1；3. 借阅记录列表与图书库存同步更新 \\
\hline
\textbf{基本事件流} & 1. 管理员进入借阅操作页面2. 管理员录入图书 ISBN、借阅人姓名，提交借阅请求3. 系统根据 ISBN 查询目标图书，校验图书是否存在4. 系统校验图书库存是否大于 05. 系统校验该借阅人是否存在该图书的未归还借阅记录6. 所有校验通过，系统生成借阅记录，记录借阅时间、借阅人、图书 ISBN 等信息7. 系统扣减对应图书的库存数量8. 系统提示 ，更新借阅记录列表，完成操作 \\
\hline
\textbf{意外事件流} & 1. 若 ISBN 对应的图书不存在，系统提示 ，终止操作2. 若图书库存为 0，系统提示 ，终止操作3. 若该借阅人已有该图书的未归还记录，系统提示 "无法重复借阅"，终止操作4. 若借阅人姓名为空，系统提示，终止提交5. 若数据保存失败，系统提示 ，回滚库存扣减与记录生成操作 \\
\hline
\textbf{活动图} & \\
\hline
\end{tabular}
\caption{借阅图书用例分析}
\end{table}

\paragraph{归还图书用例及建模}
\leavevmode 
\begin{table}[htbp]
\centering
\begin{tabular}{|p{3cm}|p{12cm}|}
\hline
\textbf{用例描述} & 图书管理员为借阅人办理图书归还手续，系统校验有效借阅记录后，完成归还操作，同步恢复图书库存、更新借阅记录状态 \\
\hline
\textbf{参与者} & 图书管理员 \\
\hline
\textbf{前置条件} & 1. 管理员已进入借阅管理模块；2. 系统存在该借阅人 + 对应 ISBN 的未归还借阅记录；3. 系统数据已正常加载 \\
\hline
\textbf{后置条件} & 1. 对应借阅记录更新归还时间，状态标记为 “已归还”；2. 对应图书的库存数量增加 1；3. 借阅记录列表与图书库存同步更新 \\
\hline
\textbf{基本事件流} & 1. 管理员，进入归还操作页面2. 管理员录入图书 ISBN、借阅人姓名，提交归还请求3. 系统查询匹配的、未归还的借阅记录，校验记录是否存在4. 系统根据 ISBN 查询目标图书，校验图书是否存在5. 校验通过，系统更新借阅记录的归还时间与状态6. 系统恢复对应图书的库存数量7. 系统提示 “图书归还成功”，更新借阅记录列表，完成操作 \\
\hline
\textbf{意外事件流} & 1. 若未找到匹配的未归还借阅记录，系统提示 ，终止操作2. 若 ISBN 对应的图书不存在，系统提示 ，终止操作3. 若数据更新失败，系统提示 ，回滚库存恢复与记录更新操作 \\
\hline
\textbf{活动图} & \\
\hline
\end{tabular}
\caption{归还图书用例分析}
\end{table}

\paragraph{查看借阅记录用例及建模}
\leavevmode 
\begin{table}[htbp]
\centering
\begin{tabular}{|p{3cm}|p{12cm}|}
\hline
\textbf{用例描述} & 用户选择查看维度，系统筛选并展示对应的借阅记录，支持多场景的记录追溯 \\
\hline
\textbf{参与者} & 图书管理员、普通用户 \\
\hline
\textbf{前置条件} & 1. 用户已进入借阅管理模块；2. 系统借阅记录与图书数据已正常加载 \\
\hline
\textbf{后置条件} & 系统展示匹配的借阅记录列表，无匹配结果时给出明确提示，不修改存储数据 \\
\hline
\textbf{基本事件流} & 1. 用户，进入记录查询页面2. 用户选择查看维度3. 用户补充筛选条件，提交查询请求4. 系统按筛选条件遍历借阅记录，匹配对应数据5. 系统为每条记录匹配对应的图书书名，渲染借阅记录列表，完成操作 \\
\hline
\textbf{意外事件流} & 1. 若无匹配的借阅记录，系统提示 ，展示空列表2. 若选择指定借阅人查询但未输入姓名，系统提示 ，终止查询3. 若数据未加载完成，系统提示 ，终止查询 \\
\hline
\textbf{活动图} & \\
\hline
\end{tabular}
\caption{查看借阅记录用例分析}
\end{table}

\subsection{面向数据流建模}
\subsubsection{数据流图}
数据流图如3.2.1所示
\begin{center}
图3.2.1数据流图
\end{center}

\subsubsection{控制流图}
从控制方向展现。控制流图如3.2.2所示
\begin{center}
图3.2.2控制流图
\end{center}

\subsection{面向类建模}
\subsubsection{类包分析}
包分析图如3.3.1所示
\begin{center}
图3.3.1包分析图
\end{center}

\subsection{行为建模}
\subsubsection{状态图}
该系统与书籍密不可分，书籍的状态十分重要，状态图如3.4.1所示。
\begin{center}
图3.4.1状态图
\end{center}

\subsubsection{时序图}
时序图如3.4.2所示。
\begin{center}
图3.4.2时序图
\end{center}

\section{技术环境}
\subsection{硬件环境}
Intel Core i3 或同等性能处理器，内存4GB 及以上，有100MB 及以上可用空间。

\subsection{软件环境}
操作系统为Windows 7/10/11，macOS，Linux，Python版本Python 3.6 及以上。

\end{document}